\documentclass[10pt,a4paper]{amsart}
\usepackage{amsmath,amssymb,amsthm,graphicx}
% Retained preamble for the Stacks Project chapter divisors.tex (Divisors).
% The chapter file pinned-original-source.tex carries only \input{preamble} on its first line, so no
% preamble text lives inside the chapter itself.  The Stacks Project uses one single shared preamble for
% every chapter; the snapshot stages it verbatim in sources/stacks-chow/preamble.tex.  The block below is
% copied mechanically, line for line, from that file's lines 14-18 (the Xy-pic setup) and lines 154-217
% (the native theorem declarations together with the mathematical macros).  Nothing was added, removed or
% redefined.  The wrapper class and the amsmath/amssymb/amsthm/graphicx base packages are supplied by the
% assembler, so the class/\externaldocument part of the shared preamble is not repeated here.
% For commutative diagrams we use Xy-pic
\usepackage[all]{xy}

% We use 2cell for 2-commutative diagrams.
\xyoption{2cell}
% Theorem environments.
%
\theoremstyle{plain}
\newtheorem{theorem}[subsection]{Theorem}
\newtheorem{proposition}[subsection]{Proposition}
\newtheorem{lemma}[subsection]{Lemma}

\theoremstyle{definition}
\newtheorem{definition}[subsection]{Definition}
\newtheorem{example}[subsection]{Example}
\newtheorem{exercise}[subsection]{Exercise}
\newtheorem{situation}[subsection]{Situation}

\theoremstyle{remark}
\newtheorem{remark}[subsection]{Remark}
\newtheorem{remarks}[subsection]{Remarks}

\numberwithin{equation}{subsection}

% Macros
%
\def\lim{\mathop{\mathrm{lim}}\nolimits}
\def\colim{\mathop{\mathrm{colim}}\nolimits}
\def\Spec{\mathop{\mathrm{Spec}}}
\def\Hom{\mathop{\mathrm{Hom}}\nolimits}
\def\Ext{\mathop{\mathrm{Ext}}\nolimits}
\def\SheafHom{\mathop{\mathcal{H}\!\mathit{om}}\nolimits}
\def\SheafExt{\mathop{\mathcal{E}\!\mathit{xt}}\nolimits}
\def\Sch{\mathit{Sch}}
\def\Mor{\mathop{\mathrm{Mor}}\nolimits}
\def\Ob{\mathop{\mathrm{Ob}}\nolimits}
\def\Sh{\mathop{\mathit{Sh}}\nolimits}
\def\NL{\mathop{N\!L}\nolimits}
\def\CH{\mathop{\mathrm{CH}}\nolimits}
\def\proetale{{pro\text{-}\acute{e}tale}}
\def\etale{{\acute{e}tale}}
\def\QCoh{\mathit{QCoh}}
\def\Ker{\mathop{\mathrm{Ker}}}
\def\Im{\mathop{\mathrm{Im}}}
\def\Coker{\mathop{\mathrm{Coker}}}
\def\Coim{\mathop{\mathrm{Coim}}}

% Boxtimes
%
\DeclareMathSymbol{\boxtimes}{\mathbin}{AMSa}{"02}

%
% Macros for moduli stacks/spaces
%
\def\QCohstack{\mathcal{QC}\!\mathit{oh}}
\def\Cohstack{\mathcal{C}\!\mathit{oh}}
\def\Spacesstack{\mathcal{S}\!\mathit{paces}}
\def\Quotfunctor{\mathrm{Quot}}
\def\Hilbfunctor{\mathrm{Hilb}}
\def\Curvesstack{\mathcal{C}\!\mathit{urves}}
\def\Polarizedstack{\mathcal{P}\!\mathit{olarized}}
\def\Complexesstack{\mathcal{C}\!\mathit{omplexes}}
% \Pic is the operator that assigns to X its picard group, usage \Pic(X)
% \Picardstack_{X/B} denotes the Picard stack of X over B
% \Picardfunctor_{X/B} denotes the Picard functor of X over B
\def\Pic{\mathop{\mathrm{Pic}}\nolimits}
\def\Picardstack{\mathcal{P}\!\mathit{ic}}
\def\Picardfunctor{\mathrm{Pic}}
\def\Deformationcategory{\mathcal{D}\!\mathit{ef}}

% --- appended 2026-10-12 (Chief): the project own definitions of the reference/slogan/history
% display environments, copied mechanically from the shared Stacks preamble lines 8-12
% (\usepackage{verbatim} plus the three \newenvironment lines, which the chapter preamble keeps
% verbatim from the project preamble). Nothing added, removed or redefined: these environments
% carry no mathematical content and the project itself typesets them as comments. Without these
% three lines a counted statement whose own text uses them fails with Environment slogan undefined.
\usepackage{verbatim}
\newenvironment{reference}{\comment}{\endcomment}
\newenvironment{slogan}{\comment}{\endcomment}
\newenvironment{history}{\comment}{\endcomment}
\usepackage{cleveref}
\crefname{originallemma}{original lemma}{original lemmas}
\Crefname{originallemma}{Original lemma}{Original lemmas}
\makeatletter
\newlabel{algebra-lemma-UFD-normal}{{0AFV}{source}{}{sourcepointer-1}{}}
\newlabel{algebra-lemma-UFD-normal@cref}{{[originallemma][0][]0AFV}{[1][0][]source}{}{}{}}
\makeatother
\crefname{originallemma}{original lemma}{original lemmas}
\Crefname{originallemma}{Original lemma}{Original lemmas}
\makeatletter
\newlabel{algebra-lemma-characterize-UFD-height-1}{{0AFT}{source}{}{sourcepointer-2}{}}
\newlabel{algebra-lemma-characterize-UFD-height-1@cref}{{[originallemma][0][]0AFT}{[1][0][]source}{}{}{}}
\makeatother
\crefname{originallemma}{original lemma}{original lemmas}
\Crefname{originallemma}{Original lemma}{Original lemmas}
\makeatletter
\newlabel{algebra-lemma-normal-domain-intersection-localizations-height-1}{{031T}{source}{}{sourcepointer-3}{}}
\newlabel{algebra-lemma-normal-domain-intersection-localizations-height-1@cref}{{[originallemma][0][]031T}{[1][0][]source}{}{}{}}
\makeatother
\crefname{originaldefinition}{original definition}{original definitions}
\Crefname{originaldefinition}{Original definition}{Original definitions}
\makeatletter
\newlabel{definition-invertible-sheaf-effective-Cartier-divisor}{{01WX}{source}{}{sourcepointer-4}{}}
\newlabel{definition-invertible-sheaf-effective-Cartier-divisor@cref}{{[originaldefinition][0][]01WX}{[1][0][]source}{}{}{}}
\makeatother
\crefname{originalequation}{original equation}{original equations}
\Crefname{originalequation}{Original equation}{Original equations}
\makeatletter
\newlabel{equation-c1}{{0BE7}{source}{}{sourcepointer-5}{}}
\newlabel{equation-c1@cref}{{[originalequation][0][]0BE7}{[1][0][]source}{}{}{}}
\makeatother
\crefname{originallemma}{original lemma}{original lemmas}
\Crefname{originallemma}{Original lemma}{Original lemmas}
\makeatletter
\newlabel{lemma-normal-c1-injective}{{0BE8}{source}{}{sourcepointer-6}{}}
\newlabel{lemma-normal-c1-injective@cref}{{[originallemma][0][]0BE8}{[1][0][]source}{}{}{}}
\makeatother
\crefname{originallemma}{original lemma}{original lemmas}
\Crefname{originallemma}{Original lemma}{Original lemmas}
\makeatletter
\newlabel{lemma-weil-divisor-is-cartier-UFD}{{0AGA}{source}{}{sourcepointer-7}{}}
\newlabel{lemma-weil-divisor-is-cartier-UFD@cref}{{[originallemma][0][]0AGA}{[1][0][]source}{}{}{}}
\makeatother
\title{For a locally Noetherian integral scheme, UFD local rings are equivalent to normality together with surjectivity of the Picard to class group map}
\author{The Stacks Project Authors}
\date{stacks-project:divisors.tex:local-snapshot-20261012; GNU FDL 1.2 or later}
\begin{document}
\maketitle
\noindent\textit{Editorial note.} All delivered mathematical text is the author's own, copied from the pinned Stacks Project chapter file pinned-original-source.tex; no statement or proof step was rewritten, reordered or removed.\par\medskip
\noindent\textit{Editorial note.} Following the multi-claim convention for this source, source\_url is the source-level URL https://stacks.math.columbia.edu/ and the source label lemma-local-rings-UFD-c1-bijective is carried in claim\_key and here rather than in a per-claim tag URL.\par\medskip
\noindent\textit{Editorial note.} The preamble is the chapter's shared Stacks Project preamble, retained as sources/stacks-divisors/retained-preamble.tex: a line-for-line excerpt of the shared Stacks preamble staged with this snapshot. No macro was invented or redefined, and this card adds no preamble text of its own.\par\medskip
\noindent\textit{Editorial note.} Statement span 6907-6918 and proof span 6920-6953 come from the halfA pool card operations/pool-cards/stacks-divisors-halfA.json and were re-read in the pinned chapter file before this card was written.\par\medskip
\noindent\textit{Editorial note.} The proof refers to the class group map by the author's equation label and to standard definitions of effective Cartier divisors and their sections; the references are part of the counted span and are handled as navigation pointers.\par\medskip
\noindent\textit{Editorial note.} The second direction uses the same argument twice, once to place an element in the local ring and once to show that every element of the prime is a multiple of it; the author's phrasing is preserved.\par\medskip
\noindent\textit{Editorial note.} 7 cross-references occur inside the counted spans; they are listed in external\_references as navigation pointers only.\par\medskip
\noindent\textit{Allowed standard prerequisites (editorial).}\begin{itemize}
\item A unique factorisation domain is normal, and the class group map of a normal integral scheme is injective, both cited from the algebra chapter and the same chapter. Cited at proof lines 6920-6923. Gives normality and injectivity in the first direction.
\item A prime divisor on a scheme whose local rings are unique factorisation domains is an effective Cartier divisor, cited from the same chapter. Cited at proof lines 6924-6926. Realizes every prime divisor as a divisor of a section, which gives surjectivity of the class group map.
\item An element of the fraction field of a normal Noetherian domain with prescribed orders at height one primes, with order one at a chosen one and order zero elsewhere, lies in the domain, cited from the algebra chapter. Cited at proof lines 6944-6947. Produces a principal generator of each height one prime in the second direction.
\item A Noetherian normal domain all of whose height one primes are principal is a unique factorisation domain, cited from the algebra chapter. Cited at proof lines 6949-6952. Concludes that the local rings are unique factorisation domains.
\end{itemize}\medskip
\ifdefined\hypertarget\hypertarget{sourcepointer-1}{}\fi
\noindent\textit{Source reference.} Original lemma 0AFV: focused-seeds/AGP/sources/stacks-snapshot-20261011/tags.txt, source lines 13173--13173. Same-chapter reference used in the retained author text at source line(s) 6922. It is supplied only as a navigation pointer so that the retained text compiles; the counted item delivers the author's own statement and the author's complete proof, and no part of the referenced item is copied in or substituted for the counted argument.\par\medskip
\ifdefined\hypertarget\hypertarget{sourcepointer-2}{}\fi
\noindent\textit{Source reference.} Original lemma 0AFT: focused-seeds/AGP/sources/stacks-snapshot-20261011/tags.txt, source lines 13171--13171. Same-chapter reference used in the retained author text at source line(s) 6952. It is supplied only as a navigation pointer so that the retained text compiles; the counted item delivers the author's own statement and the author's complete proof, and no part of the referenced item is copied in or substituted for the counted argument.\par\medskip
\ifdefined\hypertarget\hypertarget{sourcepointer-3}{}\fi
\noindent\textit{Source reference.} Original lemma 031T: focused-seeds/AGP/sources/stacks-snapshot-20261011/tags.txt, source lines 3996--3996. Same-chapter reference used in the retained author text at source line(s) 3417, 6903, 6948, 7032. It is supplied only as a navigation pointer so that the retained text compiles; the counted item delivers the author's own statement and the author's complete proof, and no part of the referenced item is copied in or substituted for the counted argument.\par\medskip
\ifdefined\hypertarget\hypertarget{sourcepointer-4}{}\fi
\noindent\textit{Source reference.} Original definition 01WX: focused-seeds/AGP/sources/stacks-snapshot-20261011/tags.txt, source lines 2498--2498. Same-chapter reference used in the retained author text at source line(s) 3576, 6928, 9257. It is supplied only as a navigation pointer so that the retained text compiles; the counted item delivers the author's own statement and the author's complete proof, and no part of the referenced item is copied in or substituted for the counted argument.\par\medskip
\ifdefined\hypertarget\hypertarget{sourcepointer-5}{}\fi
\noindent\textit{Source reference.} Original equation 0BE7: focused-seeds/AGP/sources/stacks-snapshot-20261011/tags.txt, source lines 14340--14340. Same-chapter reference used in the retained author text at source line(s) 6876, 6910, 6923, 6930, 6935, 7392. It is supplied only as a navigation pointer so that the retained text compiles; the counted item delivers the author's own statement and the author's complete proof, and no part of the referenced item is copied in or substituted for the counted argument.\par\medskip
\ifdefined\hypertarget\hypertarget{sourcepointer-6}{}\fi
\noindent\textit{Source reference.} Original lemma 0BE8: focused-seeds/AGP/sources/stacks-snapshot-20261011/tags.txt, source lines 14341--14341. Same-chapter reference used in the retained author text at source line(s) 6924, 7395. It is supplied only as a navigation pointer so that the retained text compiles; the counted item delivers the author's own statement and the author's complete proof, and no part of the referenced item is copied in or substituted for the counted argument.\par\medskip
\ifdefined\hypertarget\hypertarget{sourcepointer-7}{}\fi
\noindent\textit{Source reference.} Original lemma 0AGA: focused-seeds/AGP/sources/stacks-snapshot-20261011/tags.txt, source lines 13188--13188. Same-chapter reference used in the retained author text at source line(s) 3153, 3520, 6926. It is supplied only as a navigation pointer so that the retained text compiles; the counted item delivers the author's own statement and the author's complete proof, and no part of the referenced item is copied in or substituted for the counted argument.\par\medskip
\begin{lemma}
\label{lemma-local-rings-UFD-c1-bijective}
Let $X$ be a locally Noetherian integral scheme. Consider the map
(\ref{equation-c1}) $\Pic(X) \to \text{Cl}(X)$.
The following are equivalent
\begin{enumerate}
\item the local rings of $X$ are UFDs, and
\item $X$ is normal and $\Pic(X) \to \text{Cl}(X)$
is surjective.
\end{enumerate}
In this case $\Pic(X) \to \text{Cl}(X)$ is an isomorphism.
\end{lemma}

\begin{proof}
If (1) holds, then $X$ is normal by
Algebra, Lemma \ref{algebra-lemma-UFD-normal}.
Hence the map (\ref{equation-c1}) is injective by
Lemma \ref{lemma-normal-c1-injective}. Moreover,
every prime divisor $D \subset X$ is an effective
Cartier divisor by Lemma \ref{lemma-weil-divisor-is-cartier-UFD}.
In this case the canonical section $1_D$ of $\mathcal{O}_X(D)$
(Definition \ref{definition-invertible-sheaf-effective-Cartier-divisor})
vanishes exactly along $D$ and we see that the class of $D$ is the
image of $\mathcal{O}_X(D)$ under the map (\ref{equation-c1}).
Thus the map is surjective as well.

\medskip\noindent
Assume (2) holds. Pick a prime divisor $D \subset X$.
Since (\ref{equation-c1}) is surjective there exists an invertible
sheaf $\mathcal{L}$, a regular meromorphic section $s$, and $f \in R(X)^*$
such that $\text{div}_\mathcal{L}(s) + \text{div}(f) = [D]$.
In other words, $\text{div}_\mathcal{L}(fs) = [D]$.
Let $x \in X$ and let $A = \mathcal{O}_{X, x}$. Thus $A$ is
a Noetherian local normal domain with fraction field $K = R(X)$.
Every height $1$ prime of $A$ corresponds to a prime divisor on $X$
and every invertible $\mathcal{O}_X$-module restricts to the
trivial invertible module on $\Spec(A)$. It follows that for every
height $1$ prime $\mathfrak p \subset A$ there exists an element $f \in K$
such that $\text{ord}_{A_\mathfrak p}(f) = 1$ and
$\text{ord}_{A_{\mathfrak p'}}(f) = 0$ for every other
height one prime $\mathfrak p'$. Then $f \in A$ by Algebra, Lemma
\ref{algebra-lemma-normal-domain-intersection-localizations-height-1}.
Arguing in the same fashion we see that every element $g \in \mathfrak p$
is of the form $g = af$ for some $a \in A$. Thus we see that every
height one prime ideal of $A$ is principal and $A$ is a UFD
by Algebra, Lemma \ref{algebra-lemma-characterize-UFD-height-1}.
\end{proof}

\end{document}
